\documentclass[11pt]{article}
\usepackage[margin=25mm]{geometry}
\usepackage{fontspec,unicode-math}
\setmainfont{Latin Modern Roman}
\setmathfont{Latin Modern Math}
\usepackage{amsmath,amssymb,mathtools}
\usepackage{xurl,hyperref}
\hypersetup{colorlinks=true,urlcolor=blue,linkcolor=blue}
\setcounter{secnumdepth}{0}
\setlength{\parindent}{0pt}
\setlength{\parskip}{5pt}
\setlength{\emergencystretch}{3em}
\newcommand{\sref}[2]{\hyperlink{source:#1:#2}{[S#2]}}
\newcommand{\cataloguescope}[1]{\paragraph{Source question.}#1}
\newcommand{\floor}[1]{\left\lfloor#1\right\rfloor}
\begin{document}
\section{Geometric Explanations and Polynomial Pivot Rules for Simplex}\label{statement}

\subsection{Definitions and mathematical statement}

\paragraph{Shared notation.}Standard number systems are denoted by $\mathbb N,\mathbb Z,\mathbb Q,\mathbb R,\mathbb C$; index ranges, computational models and specialized notation follow the cited problem source and the definitions below.


Polynomial time means a number of elementary steps bounded by a polynomial in the encoded input length. A decision problem is in $\mathrm{NP}$ if yes-instances have polynomial-length certificates verifiable in polynomial time; NP-hardness is under polynomial-time reductions.

A linear program optimizes $c^Tx$ over a polyhedron $\{x:Ax\leq b\}$. In the simplex method a pivot moves between adjacent basic feasible solutions; a pivot rule chooses the next such move.
\paragraph{Statement recorded in the supplied source.}
For the simplex method for linear programming, can one explain its good practical performance by a geometric property of real-world linear programs, and does there exist a polynomial-time pivot rule? \sref{statement}{2}
\paragraph{Scope qualification.}
The practical-explanation subquestion supplies no dataset or formal criterion. The pivot-rule question retains the source\textquotesingle s permitted rule and arithmetic/bit-complexity conventions.
\paragraph{Source status.}
The supplied catalogue labels this record open. Its source-status metadata was checked against the supplied snapshot on 2026-10-11. This import does not independently establish current open status or mathematical correctness. The catalogue literature assessment is dated 2026-08-21.
\subsection{Short English statement}
Explain the cited practical simplex behavior geometrically and find a polynomial-time pivot rule.
\subsection{Sources}
\begin{itemize}
\item[\hypertarget{source:statement:1}{S1}] Ulam AI and the UnsolvedMath Contributors, UnsolvedMath: A Curated Collection of Open Mathematics Problems, catalogue entry. Catalogue curation and metadata: CC BY 4.0. The supplied source question is retained; definitions and the English summary are editorial additions.
\url{https://huggingface.co/datasets/ulamai/UnsolvedMath}.
\item[\hypertarget{source:statement:2}{S2}] Convex Geometry and its Applications (2019), Oberwolfach Reports, DOI 10.4171/owr/2018/54, pp. 3217--3281
\url{https://doi.org/10.4171/owr/2018/54}.
\item[\hypertarget{source:statement:3}{S3}] Recent simplex lower bounds
\url{https://arxiv.org/abs/2403.04886}.
\textit{Underlying Oberwolfach Reports material retains its source-specific licence, including CC BY-SA 4.0 where applicable; see THIRD\_PARTY\_NOTICES.md.}
\end{itemize}
\paragraph{Preparation and review.}
Original question attributed to the cited source and UnsolvedMath contributors. Codex (GPT-6; exact serving revision not exposed) assisted with definitions, English summaries, source inspection and formatting. No new human mathematical authorship or independent proof verification is asserted. Catalogue identity checked against the complete supplied registry and existing repository statements; known flagged repetitions and obvious identical conjectures excluded. Full mathematical-equivalence/literature review has not been performed. Statement preparation and final review followed contribution prompts C and D; automated checks verify structure and identity, not open status.
\label{end:statement}
\end{document}
